% Class-compatibility preamble retained for reader excerpts of
%   Hiraku Nakajima, "Quiver Varieties and Branching", SIGMA 5 (2009) 003,
%   DOI 10.3842/SIGMA.2009.003, CC BY 4.0.
%
% The source pins the publisher class sigma, which is not part of the delivery
% baseline, so this file supplies the amsart baseline and then copies the author's
% own declaration block mechanically and unchanged (pinned-original-source.tex
% lines 43-239): the aenume environment, the Theorem/Corollary/Lemma/Proposition and
% definition-style newtheorem declarations, the equation numbering, the author's
% reference macros (\lemref and the like) and the author's mathematical macros.
% Everything from pinned source line 240 onwards is left out: those lines define
% \rf, which uses the publisher class internal dimension \y@boxdim, and the two
% colour rules that serve that one macro; nothing in the delivered excerpts uses
% \rf, \graysquare or \halfgray.
% No mathematical macro, statement, symbol or assumption is introduced, changed or
% removed by this file.  The author's package mathdots (pinned source line 20) is not
% reproduced because it is not installed in the delivery TeX installation and no
% delivered excerpt uses its commands; the author's remaining packages are kept.
\documentclass[10pt,a4paper]{amsart}
\usepackage{amssymb}
\usepackage{amsbsy,amscd,mathtools}
\usepackage{mathrsfs}
\usepackage{nicefrac}
\usepackage[all]{xy}
\usepackage{color}
\usepackage{hyperref}
\newenvironment{aenume}{%
  \begin{enumerate}%
  \renewcommand{\theenumi}{\alph{enumi}}%
  \renewcommand{\labelenumi}{\textup{(\theenumi)}}%
  }{\end{enumerate}}
%

\usepackage{verbatim}
\usepackage{version}
\usepackage{color}
\newenvironment{NB}{
\color{red}{\bf NB}. \footnotesize
}{}
% For a final version, NB disappears.
\excludeversion{NB}
% Older version
%\renewenvironment{NB}{%
%  \comment
%  }{\endcomment}
% The following is used inside NB
\newenvironment{NB2}{
\color{blue}{\bf NB}. \footnotesize
}{}
%       Theorem environments
%% \theoremstyle{plain} %% This is the default
\newtheorem{Theorem}[equation]{Theorem}
\newtheorem{Corollary}[equation]{Corollary}
\newtheorem{Lemma}[equation]{Lemma}
\newtheorem{Proposition}[equation]{Proposition}
%\newtheorem{Axiom}{Axiom}

\theoremstyle{definition}
\newtheorem{Definition}[equation]{Definition}
\newtheorem{Example}[equation]{Example}
\newtheorem{CounterExample}[equation]{Counter-Example}
\newtheorem{Notation}[equation]{Notation}
\newtheorem{Convention}[equation]{Convention}
\newtheorem{Conjecture}[equation]{Conjecture}

%\theoremstyle{remark}
\newtheorem{Remark}[equation]{Remark}
\newtheorem{Remarks}[equation]{Remarks}
%\newtheorem*{Claim}{Claim}
%\renewcommand{\theClaim}{} % to make the claim environment unnumbered
\newtheorem{Question}[equation]{Question}
%\renewcommand{\theQuestion}{}  % to make the notation environment
                               % unnumbered
\newtheorem{Exercise}[equation]{Exercise}
%\renewcommand{\theExercise}{}

\numberwithin{equation}{section}
%\numberwithin{Theorem}{section}
%\numberwithin{figure}{section}

\newcommand{\thmref}[1]{Theorem~\ref{#1}}
\newcommand{\secref}[1]{Section~\ref{#1}}
\newcommand{\lemref}[1]{Lemma~\ref{#1}}
\newcommand{\propref}[1]{Proposition~\ref{#1}}
\newcommand{\corref}[1]{Corollary~\ref{#1}}
\newcommand{\subsecref}[1]{Section~\ref{#1}}
\newcommand{\chapref}[1]{Chapter~\ref{#1}}
\newcommand{\defref}[1]{Definition~\ref{#1}}
\newcommand{\remref}[1]{Remark~\ref{#1}}
\newcommand{\remsref}[1]{Remarks~\ref{#1}}
%
%       Math definitions
%\newcommand{\lsp}[2]{{}^{#1}{#2}}
\newcommand{\lsp}[2]{{\mskip-.3mu}{}^{#1}\mskip-1mu{#2}}
\newcommand{\defeq}{\overset{\operatorname{\scriptstyle def}}{=}}
\newcommand{\interior}{\rule{6pt}{.4pt}\rule{.4pt}{6pt}} % interior
                                                         % product
\newcommand{\C}{{\mathbb C}}
\newcommand{\Z}{{\mathbb Z}}
\newcommand{\Q}{{\mathbb Q}}
\newcommand{\R}{{\mathbb R}}
% Projective spaces
\newcommand{\proj}{{\mathbb P}}
%\newcommand{\CP}{\operatorname{\C P}}
\newcommand{\CP}{\proj}
\newcommand{\RP}{\operatorname{\R P}}
% Lie Groups
\newcommand{\SL}{\operatorname{\rm SL}}
\newcommand{\SU}{\operatorname{\rm SU}}
\newcommand{\GL}{\operatorname{GL}}
\newcommand{\PGL}{\operatorname{PGL}}
\newcommand{\U}{\operatorname{\rm U}}
\newcommand{\SO}{\operatorname{\rm SO}}
\newcommand{\grpSp}{\operatorname{\rm Sp}}
% Lie algebras
\newcommand{\algsl}{\operatorname{\mathfrak{sl}}} % because \sl="slant"
\newcommand{\su}{\operatorname{\mathfrak{su}}}
\newcommand{\gl}{\operatorname{\mathfrak{gl}}}
\newcommand{\algu}{\operatorname{\mathfrak{u}}}
\newcommand{\so}{\operatorname{\mathfrak{so}}}
\newcommand{\algsp}{\operatorname{\mathfrak{sp}}}
\newcommand{\g}{{\mathfrak g}}
\newcommand{\h}{{\mathfrak h}}
% Useful symbols
\newcommand{\Spec}{\operatorname{Spec}\nolimits}
\newcommand{\Proj}{\operatorname{Proj}\nolimits}
\newcommand{\End}{\operatorname{End}}
\newcommand{\Hom}{\operatorname{Hom}}
\newcommand{\Ext}{\operatorname{Ext}}
\newcommand{\Ker}{\operatorname{Ker}}
\newcommand{\Coker}{\operatorname{Coker}}
\newcommand{\Ima}{\operatorname{Im}}
\newcommand{\coker}{\operatorname{coker}}
\newcommand{\Endskew}{\operatorname{Endskew}}
\newcommand{\im}{\mathop{\text{\rm im}}\nolimits}
%\newcommand{\dim}{\mathop{\text{\rm dim}}\nolimits}
\newcommand{\rank}{\operatorname{rank}}
\newcommand{\ind}{\mathop{\text{\rm ind}}\nolimits}
\newcommand{\tr}{\operatorname{tr}}
\newcommand{\Tr}{\operatorname{Tr}}
\newcommand{\DB}{\overline{\partial}}
\newcommand{\pd}[2]{\frac{\partial#1}{\partial#2}}
\newcommand{\deriv}[1]{\frac{d}{d#1}}
\newcommand{\id}{\operatorname{id}}
\newcommand{\ve}{\varepsilon}
%
% notations used in this paper
%
% \newcommand{\vin}[1]{\operatorname{in}(#1)} % incoming vertex
% \newcommand{\vout}[1]{\operatorname{out}(#1)} % outgoing vertex
\newcommand{\vin}[1]{\operatorname{i}(#1)} % incoming vertex
\newcommand{\vout}[1]{\operatorname{o}(#1)} % outgoing vertex
\newcommand{\bM}{{\mathbf M}} % the vector space
\newcommand{\bN}{{\mathbf N}} % the vector space
\newcommand{\M}{{\mathfrak M}} % the moduli space
\newcommand{\Mreg}{\M^{\operatorname{s}}}
\newcommand{\La}{{\mathfrak L}} % the Lagrangian variety
\newcommand{\Fi}{{\mathfrak F}} % the fixed point variety
\newcommand{\dslash}{/\!\!/} % for algebro-geometric quotient (double slash)
\newcommand{\bv}{{\mathbf v}} % the vector v
\newcommand{\bw}{{\mathbf w}} % the vector w
\newcommand{\bC}{{\mathbf C}} % the Cartan matrix
\newcommand{\bA}{{\mathbf A}} % the adjacency matrix
\newcommand{\bI}{{\mathbf I}} % the identity matrix
\newcommand{\bn}{{\mathbf n}} % the kernel of the extended Cartan matrix
\newcommand{\bu}{{\mathbf u}}
\newcommand{\be}{{\mathbf e}}
\newcommand{\bq}{{\mathbf q}}
\newcommand{\bs}{\mathbf s}
\newcommand{\bc}{\mathbf c}
\newcommand{\codim}{\operatorname{codim}} % codimension
\newcommand{\gr}{\operatorname{gr}} % grade
\newcommand{\topdeg}{{\operatorname{top}}} % top degree
\newcommand{\Uq}{{\mathbf U}_q(\mathfrak g)} % the QUE algebra
\newcommand{\Un}{{\mathbf U}_q}
\newcommand{\Um}{\Tilde{\mathbf U}_q(\mathfrak g)} % the modified
                                                 % enveloping algebra
\newcommand{\Ua}{{\mathbf U}_q(\widehat{\mathfrak g})} % the quantum affine
                                                 % algebra
\newcommand{\Uae}{{\mathbf U}_{\varepsilon}(\widehat{\mathfrak g})} %
                                                 % the quantum affine
                                                 % algebra
\newcommand{\Ul}{{\mathbf U}_q({\mathbf L}{\mathfrak g})} % the quantum
                                                      % loop algebra
\newcommand{\Uma}{\Tilde{\mathbf U}^\bw_q(\widehat{\mathfrak g})}
%                                                 the quantum affine algebra
\newcommand{\Ui}{{\mathbf U}^{\Z}_q(\mathfrak g)}
\newcommand{\Ue}{{\mathbf U}_{\varepsilon}(\mathfrak g)}
\newcommand{\Ueone}{{\mathbf U}_{\varepsilon = 1}(\mathfrak g)}
\newcommand{\Uli}{{\mathbf U}^{\Z}_q({\mathbf L}\mathfrak g)}
\newcommand{\Ule}{{\mathbf U}_{\varepsilon}({\mathbf L}\mathfrak g)}
\newcommand{\bU}{\mathbf U} % the enveloping algebra
\newcommand{\Uqm}{{\widetilde{\mathbf U}_q}} % its modified form
\newcommand{\Pa}{{\mathfrak P}} % Hecke correspondence
\newcommand{\HomE}{\operatorname{E}}
\newcommand{\HomL}{\operatorname{L}}
\newcommand{\Zreg}{Z^{\operatorname{reg}\nolimits}}
\newcommand{\hiw}{{L}} % the highest weight module
%      \interval is used to provide better spacing after a [ that
%      is used as a closing delimiter.
\newcommand{\interval}[1]{\mathinner{#1}}
\newcommand{\shfO}{\mathcal O}
\newcommand{\Mp}{\widetilde\M}
\newcommand{\Mpp}{\widehat\M}
\newcommand{\Wedge}{{\textstyle \bigwedge}}
%
\newcommand{\Zw}{Z(\bw)}
\newcommand{\Mw}{\M(\bw)}
\newcommand{\Law}{\La(\bw)}
%\newcommand{\Zw}{Z(\bullet,\bullet;\bw)}
%\newcommand{\Mw}{\M(\bullet,\bw)}
%\newcommand{\Law}{\La(\bullet,\bw)}
\newcommand{\ngpiso}{\overset{.}{=}}
\newcommand{\idl}{\mathfrak I}
\newcommand{\Ktop}{K_{\operatorname{top}}}
\newcommand{\Kztop}{K_{0,\operatorname{top}}}
\newcommand{\Kotop}{K_{1,\operatorname{top}}}

\newcommand{\qH}{\mathbb H}
\newcommand{\D}{\mathcal D}
\newcommand{\A}{\mathcal A}
\newcommand{\incarrow}{\xrightarrow{\mathrm{inclusion}}}
\newcommand{\proarrow}{\xrightarrow{\mathrm{projection}}}
\usepackage{cleveref}
\crefname{originaltheorem}{original theorem}{original theorems}
\Crefname{originaltheorem}{Original theorem}{Original theorems}
\makeatletter
\newlabel{thm:affine}{{2.6}{source}{}{sourcepointer-1}{}}
\newlabel{thm:affine@cref}{{[originaltheorem][0][]2.6}{[1][0][]source}{}{}{}}
\makeatother
\crefname{originalequation}{original equation}{original equations}
\Crefname{originalequation}{Original equation}{Original equations}
\makeatletter
\newlabel{eq:trivial}{{1.9}{source}{}{sourcepointer-2}{}}
\newlabel{eq:trivial@cref}{{[originalequation][0][]1.9}{[1][0][]source}{}{}{}}
\makeatother
\crefname{originalequation}{original equation}{original equations}
\Crefname{originalequation}{Original equation}{Original equations}
\makeatletter
\newlabel{eq:diagram}{{2.14}{source}{}{sourcepointer-3}{}}
\newlabel{eq:diagram@cref}{{[originalequation][0][]2.14}{[1][0][]source}{}{}{}}
\makeatother
\title{The fibre-dimension bound and the semismallness of the projective morphism between two quiver varieties}
\author{Hiraku Nakajima}
\date{doi:10.3842/SIGMA.2009.003; CC BY 4.0}
\begin{document}
\maketitle
\noindent\textit{Editorial note.} All delivered mathematical text is the author's own, copied from the pinned publisher source sources/algebra-fl01-quiver-varieties-and-branching/pinned-original-source.tex; no statement or proof step was rewritten, reordered or removed.\par\medskip
\noindent\textit{Editorial note.} Preamble provenance: as for the other cards of this source, the retained preamble supplies the amsart baseline in place of the publisher class sigma and copies the author's declaration block of pinned source lines 43-239 unchanged.\par\medskip
\noindent\textit{Editorial note.} Bibliography provenance: the statement carries the author's own attribution cf. to two references (Na-quiver section 6.11 and Na-alg sections 10.4 and 10.11) and the proof repeats the first of them. These citations are resolved with sources/algebra-fl01-quiver-varieties-and-branching/bibliography.tex, an excerpt of the author's own thebibliography environment (pinned source lines 4098-4257) copied character for character; no entry was rewritten and no citation key was changed.\par\medskip
\noindent\textit{Editorial note.} External references are navigation pointers only; every deliberately kept step of the author's argument stays inside the counted spans. Their printed numbers are the author's own numbering under the author's newtheorem and equation-numbering declarations, computed mechanically by the lane's scanner.\par\medskip
\noindent\textit{Allowed standard prerequisites (editorial).}\begin{itemize}
\item The morphism pi from M\_zeta(V,W) to M\_\{zeta\textasciicircum{}bullet\}(V,W), the strata (M\_\{zeta\textasciicircum{}bullet\})\_(G-hat) and the chamber condition making M\_zeta(V,W) nonsingular. Set up earlier in the same subsection (source lines 1324-1552, the author's subsections on strata and on the local structure). Fixes the map, its target stratification and the smoothness assumption under which the fibre-dimension bound is stated.
\item Semismallness of a morphism with respect to a stratification. Used in part (2) of the statement (source lines 1566-1574). Names the property that part (2) establishes for the restricted map pi into its image.
\end{itemize}\medskip
\ifdefined\hypertarget\hypertarget{sourcepointer-1}{}\fi
\noindent\textit{Source reference.} Original theorem 2.6: sources/algebra-fl01-quiver-varieties-and-branching/pinned-original-source.tex, source lines 944--944. Navigation pointer only: part (1) of the counted proof invokes the author's theorem that a certain moduli space is an affine algebraic manifold, in one clause; its statement is not reproduced here and the cardinality-free content it supplies is declared as a named standard prerequisite.\par\medskip
\ifdefined\hypertarget\hypertarget{sourcepointer-2}{}\fi
\noindent\textit{Source reference.} Original equation 1.9: sources/algebra-fl01-quiver-varieties-and-branching/pinned-original-source.tex, source lines 660--660. Navigation pointer only: it labels the author's displayed trivial factor of the space bM, which the counted proof mentions in one clause; the stub keeps the author's own equation number resolvable without reproducing that display.\par\medskip
\ifdefined\hypertarget\hypertarget{sourcepointer-3}{}\fi
\noindent\textit{Source reference.} Original equation 2.14: sources/algebra-fl01-quiver-varieties-and-branching/pinned-original-source.tex, source lines 1494--1494. Navigation pointer only: it labels the author's commutative diagram for the local isomorphism, which the counted proof names before drawing its conclusion; the diagram itself is not reproduced.\par\medskip
\begin{Proposition}[cf.~\protect{\cite[6.11]{Na-quiver},
\cite[10.4, 10.11]{Na-alg}}]\label{prop:semismall}
Let $\zeta$, $\zeta^\bullet$ be as above and consider
$\pi\colon\M_\zeta(V,\linebreak[3]W)\to \M_{\zeta^\bullet}(V,W)$.
%
  We further assume that $\zeta$ is in a chamber so that
  $\M_\zeta(V,W)$ is nonsingular.

\textup{(1)} Take a point $x$ in a stratum $(\M_{\zeta^\bullet})_{(\widehat
  G)}$. If $\pi^{-1}(x)\neq \varnothing$, then we have
\begin{equation*}
  \dim \pi^{-1}(x)
  \le \tfrac12
  \big( \dim \M_\zeta(V,W) - \dim (\M_{\zeta^\bullet})_{(\widehat G)}\big).
\end{equation*}

\textup{(2)} Replace the target of $\pi$ by the image, hence consider
$\pi\colon \M_{\zeta}(V,W)\to \pi(\M_\zeta(V,W))$. It is semismall
with respect to the stratification $\pi(\M_\zeta(V,W)) = \bigsqcup
(\M_{\zeta^\bullet})_{(\widehat G)}$ where $\widehat G$ runs over the
conjugacy classes of subgroups of $G_V$ such that
$(\M_{\zeta^\bullet})_{(\widehat G)}$ is contained in $\pi(\M_\zeta(V,W))$.
\end{Proposition}

\begin{proof}
  (1) This was proved under the assumption
  $\Mreg_{\zeta^\bullet}(V,W)\neq\varnothing$ in \cite[6.11]{Na-quiver},
  but its proof actually gives the above. The point of the proof was
  that the f\/iber $\widehat\pi^{-1}(0)$ is a subvariety in
  $\widehat\M_\zeta(\Hat{V},\Hat{W})$ (divided by the trivial factor
  \eqref{eq:trivial}), which is an af\/f\/ine algebraic manifold by
  \thmref{thm:affine}.

(2) By \eqref{eq:diagram} and the subsequent remark, for each stratum
$(\M_{\zeta^\bullet})_{(\widehat G)}$, the restriction of $\pi$ to the
preimage $\pi^{-1}((\M_{\zeta^\bullet})_{(\widehat G)})$ is a f\/iber bundle.
Moreover, by (1), there exists a stratum $(\M_{\zeta^\bullet})_{(\widehat
  G)}$ such that $\dim (\M_{\zeta^\bullet})_{(\widehat G)} = \dim
\M_\zeta(V,W)$. Therefore we have $\dim \pi(\M_\zeta(V,W)) = \dim
\M_\zeta(V,W)$. Hence (1) implies the assertion.
\end{proof}

% Bibliography excerpt retained for reader excerpts of
%   Hiraku Nakajima, "Quiver Varieties and Branching", SIGMA 5 (2009) 003,
%   DOI 10.3842/SIGMA.2009.003, CC BY 4.0.
%
% Every entry below is copied character for character from the author's own
% thebibliography environment of the pinned publisher source
% (pinned-original-source.tex lines 4098-4257).  No entry was rewritten,
% reordered, merged, split or invented, and no citation key was changed;
% only the entries that the delivered excerpts cite are kept.
\begin{thebibliography}{99}

\bibitem{CB}
Crawley-Boevey W., Geometry of the moment map for representations of
  quivers, {\it Compositio Math.} \textbf{126} (2001),   257--293.
 % \MR{MR1834739 (2002g:16021)}


\bibitem{Kac}
 Kac V.G., Inf\/inite-dimensional Lie algebras, 3rd ed., Cambridge
  University Press, Cambridge, 1990. %\MR{MR1104219 (92k:17038)}


\bibitem{KS}
Kashiwara M., Saito Y., Geometric construction of crystal bases, {\it Duke
  Math. J.} \textbf{89} (1997),  9--36,  \mbox{\href{http://arxiv.org/abs/q-alg/9606009}{q-alg/9606009}}. %\MR{MR1458969 (99e:17025)}


\bibitem{Na-quiver}
Nakajima H., Instantons on {ALE} spaces, quiver varieties, and
  Kac--Moody algebras, {\it Duke Math. J.} \textbf{76} (1994), 365--416.
  %\MR{MR1302318 (95i:53051)}


\bibitem{Na-proc}
Nakajima H., Varieties associated with quivers, in Representation Theory of
  Algebras and Related Topics (Mexico City, 1994), {\it CMS Conf. Proc.}, Vol.~19,
  Amer. Math. Soc., Providence, RI, 1996, 139--157.
  %\MR{MR1388562   (97m:16022)}


\bibitem{Na-alg}
Nakajima H., Quiver varieties and Kac--Moody algebras, {\it Duke Math. J.}
  \textbf{91} (1998), 515--560. %\MR{MR1604167 (99b:17033)}


\bibitem{Lecture}
Nakajima H., Lectures on Hilbert schemes of points on surfaces, {\it University
  Lecture Series}, Vol.~18, American Mathematical Society, Providence, RI, 1999.
  %\MR{MR1711344 (2001b:14007)}


\bibitem{Na-qaff}
Nakajima H., Quiver varieties and f\/inite-dimensional representations of
  quantum af\/f\/ine algebras, {\it J. Amer. Math. Soc.} \textbf{14} (2001),
  145--238, \href{http://arxiv.org/abs/math.QA/9912158}{math.QA/9912158}. %\MR{MR1808477 (2002i:17023)}


\bibitem{Na-tensor}
Nakajima H., Quiver varieties and tensor products, {\it Invent. Math.}
  \textbf{146} (2001), 399--449, \href{http://arxiv.org/abs/math.QA/0103008}{math.QA/0103008}. %\MR{MR1865400 (2003e:17023)}


\bibitem{Na-reflect}
Nakajima H.,  Ref\/lection functors for quiver varieties and {W}eyl group
  actions, {\it Math. Ann.} \textbf{327} (2003), 671--721.
  %\MR{MR2023313   (2004k:16036)}


\bibitem{Na-ADHM}
Nakajima H.,  Sheaves on $ALE$ spaces and quiver varieties, {\it Mosc. Math. J.}
  \textbf{7} (2007),  699--722.


\bibitem{Saito}
Saito Y., Crystal bases and quiver varieties, {\it Math. Ann.} \textbf{324}
  (2002),  675--688, \href{http://arxiv.org/abs/math.QA/0111232}{math.QA/0111232}. %\MR{MR1942245 (2004a:17023)}



\end{thebibliography}
\end{document}
